% Image credits for the photographs and AI illustrations of Book 5 (University, Year 3).
% Machine-readable license records live in images/book5/CREDITS.md; keep the
% two files in sync. The Book 5 writer owns this file: add one \item per
% photograph (inside the itemize, by chapter) as photographs land.
\chapter*{Créditos de las imágenes}

Los dibujos, mapas y gráficos de este libro son originales; los mapas y gráficos
elaborados a partir de datos publicados citan su fuente en el pie de figura o en
esta lista. Las fotografías se usan bajo sus respectivas licencias libres, con
gratitud a sus autores; los enlaces completos a la fuente y las URL de las
licencias de cada fotografía figuran en \texttt{images/book5/CREDITS.md} en el
repositorio del libro.

\bigskip
Junto a las fotografías y a los dibujos originales, algunas figuras son
ilustraciones generadas por IA, producidas con la generación de imágenes de OpenAI
a partir de instrucciones escritas por los editores del libro, y revisadas por su
exactitud científica antes de incluirlas. Muestran solo escenas genéricas, nunca un
lugar real, una persona real ni un ejemplar real. La instrucción completa de cada
imagen generada figura en \texttt{images/book5/ai/PROMPTS.md} en el repositorio.

\bigskip
\noindent Fotografías y datos, por capítulo:
\begin{itemize}\small
  \item Cap.~1 (la sismología y las imágenes de la Tierra): PREM, A.~M.~Dziewonski y D.~L.~Anderson (1981), tal como figura en ObsPy; ak135, B.~L.~N.~Kennett, E.~R.~Engdahl y R.~Buland (1995).
  \item Cap.~3 (la gravedad y la figura de la Tierra): los satélites gemelos GRACE Follow-On, concepción artística, NASA/JPL-Caltech, dominio público (Wikimedia Commons). Balance de masa de los casquetes de hielo del IMBIE Team (2021), UK Polar Data Centre, Open Government Licence v3.0.
  \item Cap.~4 (el geomagnetismo y el paleomagnetismo): Campo Geomagnético Internacional de Referencia, decimocuarta generación, IAGA Division~V Working Group V-MOD, coeficientes distribuidos por NOAA NCEI.
  \item Cap.~5 (los modelos térmicos de la litosfera): edad del fondo oceánico, M.~Seton y otros (2020), EarthByte, CC BY 4.0; relieve global ETOPO1, NOAA NCEI, dominio público.
  \item Cap.~6 (la convección del manto y las fuerzas que mueven las placas): las playas levantadas de Ballyhillin, Malin Head, A.~F.~Borchert, CC BY-SA 4.0 (Wikimedia Commons); PREM, A.~M.~Dziewonski y D.~L.~Anderson (1981).
  \item Cap.~7 (los mecanismos de deformación y las zonas de cizalla): milonita con ojos, Mikenorton, CC BY-SA 4.0 (Wikimedia Commons).
  \item Cap.~9 (la tectónica extensional y de desgarre): relieve de cuencas y sierras al sur del Valle de la Muerte, Ken Lund, CC BY-SA 2.0 (Wikimedia Commons).
  \item Cap.~14 (la geoquímica de los isótopos estables): LR04 benthic $\delta^{18}$O stack, L.~E.~Lisiecki y M.~E.~Raymo (2005).
  \item Cap.~15 (la interacción agua--roca): terrazas de travertino de Palette Spring, Mammoth Hot Springs, Jim Peaco, National Park Service, dominio público (Wikimedia Commons). Constantes de equilibrio de la base de datos phreeqc.dat de PHREEQC, D.~L.~Parkhurst y C.~A.~J.~Appelo, U.S. Geological Survey, dominio público.
  \item Cap.~16 (la estratigrafía secuencial): los Book Cliffs sobre Helper, Utah, David Mackie, dominio público (Wikimedia Commons); LR04 benthic $\delta^{18}$O stack, L.~E.~Lisiecki y M.~E.~Raymo (2005).
  \item Cap.~17 (los modelos de facies): estratificación cruzada hummocky en la Formación Arisaig, Michael C.~Rygel, CC BY-SA 3.0 (Wikimedia Commons), recortada.
  \item Cap.~18 (la geomorfología cuantitativa): la garganta del Niágara y las cataratas Horseshoe, Ken Lund, CC BY-SA 2.0 (Wikimedia Commons), recortada.
  \item Cap.~19 (la dinámica de fluidos geofísicos): vientos de media zonal del NCEP--NCAR Reanalysis~1, datos proporcionados por el NOAA PSL, Boulder, Colorado, USA; World Ocean Atlas 2023, NOAA NCEI.
  \item Cap.~20 (la dinámica atmosférica): una borrasca frente a Islandia, NASA/GSFC, MODIS Rapid Response Team, Jacques Descloitres, dominio público (Wikimedia Commons); vientos de media zonal del NCEP--NCAR Reanalysis~1, datos proporcionados por el NOAA PSL, Boulder, Colorado, USA.
  \item Cap.~21 (la dinámica oceánica): una boya Argo lanzada desde un buque oceanográfico, NOAA Global Ocean Monitoring and Observing Program, dominio público (Wikimedia Commons).
  \item Cap.~22 (los modelos climáticos): superordenadores del NASA Center for Climate Simulation, Goddard Space Flight Center, 2010, dominio público (Wikimedia Commons); temperatura del aire en superficie y radiación en la cima de la atmósfera del NCEP--NCAR Reanalysis~1, datos proporcionados por el NOAA PSL, Boulder, Colorado, USA.
  \item Cap.~24 (la glaciología): un glaciar emisario del suroeste de Groenlandia, NASA / Jim Yungel, Operation IceBridge, 2014, dominio público (Wikimedia Commons); altitudes de la superficie de Groenlandia de ETOPO1 Ice Surface, NOAA NCEI; extensión del hielo marino ártico del NSIDC Sea Ice Index, versión~4 (Fetterer y otros, 2025, doi:10.7265/a98x-0f50).
  \item Cap.~25 (la teledetección y los SIG): el interferograma 1995--2001 del centro volcánico de Three Sisters, Oregon, U.S. Geological Survey (C.~Wicks y otros, 2002), dominio público (Wikimedia Commons); espectros de laboratorio de la ECOSTRESS Spectral Library, versión~1.0 (S.~K.~Meerdink, S.~J.~Hook, D.~A.~Roberts y E.~A.~Abbott, 2019; A.~M.~Baldridge y otros, 2009), reproducidos por cortesía del Jet Propulsion Laboratory, California Institute of Technology, Pasadena, California.
  \item Cap.~26 (los interiores y las superficies planetarios): Barringer Meteor Crater, Arizona, USGS/D.~Roddy, dominio público (Wikimedia Commons); Olympus Mons, mosaico del orbitador Viking, NASA/\allowbreak JPL/\allowbreak USGS, dominio público (Wikimedia Commons). Funciones de producción y de cronología de cráteres tal como figuran en el programa Craterstats de G.~G.~Michael, según G.~Neukum, B.~A.~Ivanov y W.~K.~Hartmann (2001).
  \item Cap.~27 (las atmósferas planetarias y la habitabilidad): una red de valles secos en Marte, Mars Express HRSC, ESA/DLR/FU Berlin, CC BY-SA 3.0 IGO (Wikimedia Commons); Ligeia Mare, Titan, mosaico de radar de Cassini, NASA/JPL/USGS, dominio público (Wikimedia Commons). ajustes de la zona habitable de R.~K.~Kopparapu y otros (2013); presión de vapor del metano del NIST Chemistry WebBook.
  \item Cap.~28 (los exoplanetas): HIP~65426~b fotografiado por el telescopio espacial James Webb, NASA/ESA/CSA, A~Carter (UCSC), el equipo ERS~1386 y A.~Pagan (STScI), CC BY~4.0 (Wikimedia Commons). Planetas confirmados del NASA Exoplanet Archive, tabla Planetary Systems Composite Parameters, consultada el 23 de octubre de 2026: esta investigación ha utilizado el NASA Exoplanet Archive, gestionado por el California Institute of Technology bajo contrato con la National Aeronautics and Space Administration en el marco del Exoplanet Exploration Program. Curvas masa--radio modelizadas de L.~Zeng y otros (2016, 2019).
  \item Cap.~29 (la Vía Láctea y las galaxias): la panorámica de la Vía Láctea del proyecto GigaGalaxy Zoom, ESO/S.~Brunier, CC BY~4.0 (Wikimedia Commons); la galaxia espiral M101, Agencia Espacial Europea y NASA (Hubble, heic0602; datos Hubble K.~D.~Kuntz, F.~Bresolin, J.~Trauger, J.~Mould e Y.-H.~Chu; tratamiento de imagen Davide De~Martin), CC BY~4.0 (Wikimedia Commons). Velocidades circulares de la Galaxia de A.-C.~Eilers, D.~W.~Hogg, H.-W.~Rix y M.~K.~Ness (2019); modelo de masa de E.~Pouliasis, P.~Di~Matteo y M.~Haywood (2017); cefeidas clásicas OGLE de la Gran Nube de Magallanes, I.~Soszy\'nski y otros (2015), Acta Astronomica 65, 297, distribuidas por el CDS, Strasbourg (VizieR J/AcA/65/297).
  \item Cap.~30 (de dónde vienen los elementos): la nebulosa del Cangrejo, NASA, ESA, J.~Hester y A.~Loll (Arizona State University), dominio público (Wikimedia Commons); la kilonova de GW170817 en NGC~4993, NASA y ESA, agradecimiento a A.~J.~Levan, N.~R.~Tanvir, A.~Fruchter y O.~Fox (Hubble, heic1717a), CC BY~4.0 (Wikimedia Commons). Abundancias solares de M.~Asplund, A.~M.~Amarsi y N.~Grevesse (2021); secciones eficaces de captura de neutrones de KADoNiS v0.3 (I.~Dillmann, R.~Plag, F.~K\"appeler y T.~Rauscher); periodos de semidesintegración de la IAEA Live Chart of Nuclides (ENSDF); metalicidades del catálogo de Ginebra--Copenhague reanalizadas por L.~Casagrande y otros (2011), distribuidas por el CDS, Strasbourg (VizieR J/A+A/530/A138).
\end{itemize}
